% Paper 2, §6 — Mixed precision and stability
% Anchors: kb/notes/training/mixed-precision-and-stability.md
%          kb/excerpts/{micikevicius2017-mixed-precision,kalamkar2019-bfloat16,
%                       deepseek-v3-training,torchtitan2024,muon-moonlight2025,
%                       ma2024-bitnet,scaling-laws-precision-2024}.md

\section{Mixed precision and stability}
\label{sec:mixed-precision}

The numerical-precision axis is independent of the parameter axis,
scales the same way the parameter axis does, and is now production-
stable through FP8 at frontier scale \citep[\S 3.3]{deepseek-v3}.
The lineage runs FP32 baseline $\to$ FP16 with FP32 master weights and
dynamic loss scaling \citep{micikevicius2017-mixed-precision}
$\to$ BF16 as a drop-in replacement that abolishes loss scaling on
hardware with native BF16 Tensor Cores \citep{kalamkar2019-bfloat16}
$\to$ FP8 fine-grained quantization with promotion-to-CUDA-cores
accumulation \citep[\S 3.3]{deepseek-v3} $\to$ FP4-class formats
(Blackwell-era NVFP4, BitNet 1.58 ternary)\footnote{KB note:
\texttt{kb/notes/training/mixed-precision-and-stability.md\#3-1}.} that
cut storage and compute another factor of two but whose frontier-scale
production status is contested. Each step halves storage and roughly
doubles per-Tensor-Core throughput on supporting hardware
\citep[\S 3.3]{deepseek-v3}. The thesis of this section is that the
recipe is a small set of independently-tunable choices — storage
precision, compute precision, communication precision, accumulator
precision — and the headline 2026 fact is that DeepSeek-V3's full
14.8T-token, 2-month, 2048-H800 pre-training reported zero
irrecoverable loss spikes in FP8
\citep[abstract]{deepseek-v3}.\footnote{KB excerpt:
\texttt{kb/excerpts/deepseek-v3-training.md\#abstract}.}

\subsection{The FP32 baseline and why we leave it}
\label{sec:mp-fp32}

A pure-FP32 training step stores parameters $\theta \in \R^P$,
gradients $g \in \R^P$, and Adam moments $(m, v) \in \R^P \times \R^P$
all in 32-bit floats. The total optimizer-state footprint is $4P$ FP32
words: $P$ for parameters, $P$ for gradients, $2P$ for Adam's first and
second moments. For a $70$B parameter model the optimizer state alone is
$4 \cdot 70 \cdot 10^{9} \cdot 4 \approx 1.1\,\text{TB}$ before
activations, KV cache, or DP / TP / PP duplication
\citep[\S 2]{kalamkar2019-bfloat16}. The activation footprint is
similarly painful: a single forward pass at batch $\mathsf{B}$ and
sequence length $\mathsf{S}$ stores $L \cdot \mathsf{B} \cdot
\mathsf{S} \cdot \mathsf{D}$ FP32 elements per layer of the residual
stream, scaled by a few-times-$\mathsf{D}$ activation-checkpointing
multiplier.

The compute side compounds the storage problem. NVIDIA Tensor Cores
deliver substantially higher throughput at FP16 / BF16 / FP8 than at
FP32: a single FP8 GEMM on H100 runs at $\sim 4\times$ the rate of an
FP16 GEMM on the same silicon \citep[\S 3.3]{deepseek-v3}. Distributed
all-reduce traffic also halves with each precision step, since
gradient bandwidth is the dominant inter-node cost in DP / FSDP runs
\citep[\S 2.2.4]{torchtitan2024}.\footnote{KB excerpt:
\texttt{kb/excerpts/torchtitan2024.md\#sec-2-2-4}.} The mixed-precision
question is therefore not whether to leave FP32 — every frontier-scale
LLM does — but how far down the precision stack we can go before
training diverges or downstream quality degrades.

\subsection{FP16 with master weights and dynamic loss scaling}
\label{sec:mp-fp16}

\citet{micikevicius2017-mixed-precision} introduce the canonical
mixed-precision recipe that survived two GPU generations as the default
training configuration. Three independent ingredients
\citep[\S 2--4]{micikevicius2017-mixed-precision}\footnote{KB excerpt:
\texttt{kb/excerpts/micikevicius2017-mixed-precision.md\#sec-2}.}:

\begin{enumerate}
\item Maintain a single-precision \emph{master} copy of the weights,
$\theta^* \in \R^P$ in FP32, that accumulates the optimizer update;
the FP16 working copy $\theta \in \R^P$ is recast from $\theta^*$ for
each forward and backward pass.
\item Multiply the loss by a scale factor $S$ before backprop and
divide gradients by $S$ before the master-weight update — \emph{loss
scaling}.
\item Run reduced-precision GEMMs with FP32 accumulators: inputs in
FP16, accumulator in FP32, output cast back to FP16.
\end{enumerate}

The training step has the schematic form
\begin{equation}
\begin{aligned}
y &= \mathrm{Forward}(\theta_{\text{FP16}}, x), \\
\tilde g_{\text{FP16}} &= \mathrm{Backward}(y, S \cdot \mathcal{L}), \\
g^*_{\text{FP32}} &= \mathrm{cast}(\tilde g_{\text{FP16}}) / S, \\
\theta^*_{\text{FP32}} &\leftarrow \theta^*_{\text{FP32}} - \eta \cdot \mathrm{Optim}(g^*_{\text{FP32}}), \\
\theta_{\text{FP16}} &= \mathrm{cast}(\theta^*_{\text{FP32}}),
\end{aligned}
\label{eq:amp}
\end{equation}
verbatim from \citet[\S 2]{micikevicius2017-mixed-precision}. The
\emph{master copy} prevents the cumulative error of many small updates
from drowning under FP16 quantization: each $\eta \cdot
\mathrm{Optim}(g^*)$ may be far below FP16's smallest normal value, but
because it lands in the FP32 master, it accumulates exactly until the
running sum crosses an FP16-representable threshold and is reflected
back to $\theta_{\text{FP16}}$ on the next downcast.

Loss scaling addresses the orthogonal underflow problem in the
gradients themselves. FP16's smallest normal value is
$\sim\!\!6\times 10^{-5}$, and gradient distributions during LLM
training have a long tail well below this; in plain FP16 those
gradients underflow to zero
\citep[\S 3]{micikevicius2017-mixed-precision}. Multiplying the loss by
$S$ shifts the entire gradient distribution upward in the FP16 range
without changing the chain rule:
\begin{equation}
\tilde g \;=\; \nabla(S \cdot \mathcal{L}) \;=\; S \cdot g, \qquad
g^*_{\text{FP32}} \;=\; \mathrm{cast}(\tilde g_{\text{FP16}}) / S.
\label{eq:loss-scale}
\end{equation}
Static loss scaling picks a single $S \in [2^{10}, 2^{15}]$. Dynamic
loss scaling probes the largest $S$ that does not overflow:

\begin{lstlisting}[caption={Dynamic loss scaling (Micikevicius et al.\ 2017, \S 3).}, label=lst:dynamic-ls]
S = 2**15
overflow_count = 0
for step in range(num_steps):
    loss = forward(theta_FP16, x)
    grad_FP16 = backward(S * loss)
    if any_inf_or_nan(grad_FP16):
        S = S / 2          # halve scale
        overflow_count = 0
        continue           # SKIP this step's update
    grad_FP32 = cast_to_FP32(grad_FP16) / S
    theta_master_FP32 -= eta * Optim(grad_FP32)
    theta_FP16 = cast_to_FP16(theta_master_FP32)
    overflow_count += 1
    if overflow_count >= GROW_INTERVAL:
        S = S * 2          # double scale periodically
        overflow_count = 0
\end{lstlisting}

The reduced-precision GEMM with FP32 accumulator is the compute-side
ingredient: the inner-product summation accumulates in FP32 even when
both operands are FP16, preventing drift on large reduction dimensions
$K \gtrsim 1000$ \citep[\S 4]{micikevicius2017-mixed-precision}. NVIDIA
Volta and later architectures implement this as a single Tensor Core
instruction. The combined recipe achieves FP32-equivalent accuracy
across vision and NLP tasks at the time
\citep[\S 5]{micikevicius2017-mixed-precision}, and remained the
default for $\sim$3 years until BF16 made loss scaling unnecessary.

\subsection{BF16: trading precision for range}
\label{sec:mp-bf16}

\citet{kalamkar2019-bfloat16} make the case for the Brain
Floating-Point format: 16 bits with 8 exponent bits and 7 mantissa
bits, the same exponent width as FP32. The trade is explicit
\citep[\S 2]{kalamkar2019-bfloat16}\footnote{KB excerpt:
\texttt{kb/excerpts/kalamkar2019-bfloat16.md\#sec-2}.}:

\begin{table}[h]
\centering
\small
\begin{tabular}{lrrrl}
\toprule
Format & Bits & Exponent & Mantissa & Dyn. range $\approx$ \\
\midrule
FP32 & 32 & 8  & 23 & $10^{\pm 38}$ \\
FP16 & 16 & 5  & 10 & $6\times 10^{-5}$ to $6.55\times 10^{4}$ \\
\textbf{BF16} & 16 & \textbf{8} & \textbf{7} & $10^{\pm 38}$ (FP32-like) \\
\bottomrule
\end{tabular}
\caption{The 16-bit format zoo. BF16 keeps FP32's exponent width, at
the cost of 3 bits of mantissa relative to FP16. After
\citet[\S 2]{kalamkar2019-bfloat16}.}
\label{tab:fp16-bf16-fp32}
\end{table}

The empirical claim is sharp: ``models trained with BFLOAT16 weights,
activations, gradients, and updates match those trained with FP32 in
terms of accuracy without any change in hyper-parameters''
\citep[\S 4]{kalamkar2019-bfloat16}. BF16 is therefore a
\emph{drop-in} replacement for FP32 — no recipe changes, no loss
scaling, optionally no FP32 master weights. The 7-mantissa precision
loss is absorbed because each weight is updated thousands of times
during training and accumulated low-precision noise averages out
\citep[\S 3]{kalamkar2019-bfloat16}.

\begin{intuition}
BF16's ``free lunch'' relative to FP16 is range, not precision. Range
dominates for gradients, whose magnitudes vary by many orders across
the network and across coordinates; FP16's narrow range causes the
small-gradient tail to underflow before any per-tensor recipe sees it,
necessitating loss scaling. BF16 captures the FP32-range tail directly,
trading mantissa precision (which averages out across many updates) for
exponent range (which does not). Returning to math: the master-weights
pattern in \cref{eq:amp} is required when the working precision's
representable range fails to cover the gradient distribution; under
BF16's 8-bit exponent, it does, and the FP32 master becomes optional.
\end{intuition}

By 2023 BF16 had become the universal pre-training default for LLaMA
2 and 3, Mistral, Qwen, Gemma, and OLMo, on the strength of the
drop-in property and Ampere+ native BF16 Tensor Core support
\citep[\S 2.2.4]{torchtitan2024}.

\subsection{FP8 in production training}
\label{sec:mp-fp8}

The H100 and successor architectures provide native FP8 Tensor Cores
in two formats: \emph{E4M3} (4 exponent bits, 3 mantissa, range
$\pm 448$) and \emph{E5M2} (5 exponent, 2 mantissa, range $\pm
5.7\times 10^{4}$) \citep[\S 3.3]{deepseek-v3}. NVIDIA's recommended
\emph{hybrid} recipe uses E4M3 for forward activations and weights
and E5M2 for gradients, on the principle that gradients need more
dynamic range. DeepSeek-V3's recipe \citep[\S 3.3.1, \S
3.3.2]{deepseek-v3}\footnote{KB excerpts:
\texttt{kb/excerpts/deepseek-v3-training.md\#sec-3-3-fp8} and
\texttt{\#sec-3-3-2-finegrained}.} is the most aggressive published
production FP8 training to date, and its decisions disagree with
NVIDIA's hybrid in three structural ways.

\paragraph{Fine-grained quantization granularity.} For sub-FP16
formats the per-tensor or per-block scale $Z$ becomes a first-class
object. The FP8 GEMM $Y = X W$ is implemented as
\begin{equation}
Y = (Z_X \cdot X_{\text{FP8}}) (Z_W \cdot W_{\text{FP8}})
  = Z_X Z_W \, (X_{\text{FP8}} \, W_{\text{FP8}})_{\text{FP8 accum}},
\label{eq:fp8-gemm}
\end{equation}
with the inner GEMM on FP8 Tensor Cores and the scalar product
$Z_X Z_W$ folded in afterward \citep[\S 3.3.2, eq.\ 2]{deepseek-v3}.
The granularity at which $(Z_X, Z_W)$ are chosen is the load-bearing
choice. Three operating points, increasing in cost-per-element of
metadata and decreasing in quantization error:
\begin{itemize}
\item \emph{Per-tensor}: a single scale per Linear input or weight
matrix. TorchTitan's \texttt{torchao.float8} (dynamic, delayed, or
static scaling modes) applies this scheme selectively on
\texttt{Linear} layers \citep[\S 2.2.4]{torchtitan2024}.
\item \emph{Per-row / per-token}: one scale per row of the activation
matrix (i.e., per-token); used for GEMM input activations.
\item \emph{Per-block / per-tile}: DeepSeek-V3's choice, with $1
\times 128$ tiles for activations (per-token-per-128-channel) and $128
\times 128$ blocks for weights, with per-group scales along the inner
GEMM dimension \citep[\S 3.3.2]{deepseek-v3}.
\end{itemize}

\paragraph{All-E4M3 instead of E4M3/E5M2 hybrid.} DeepSeek-V3 uses
E4M3 \emph{everywhere} (forward, activation backward, weight backward)
and recovers the dynamic range that hybrid would have provided through
fine-grained scaling \citep[\S 3.3.2]{deepseek-v3}.

\paragraph{Promotion to CUDA cores at $N_C = 128$.} H100's FP8 Tensor
Core accumulates internally in ``around 14 bits'' rather than full
FP32, giving up to $\sim$2\% relative error on $K = 4096$ dot products
\citep[\S 3.3.2]{deepseek-v3}. The fix is to copy partial sums to
FP32 registers and accumulate in CUDA cores every $N_C = 128$
elements of the inner reduction, recovering FP32 accumulator semantics
at a small bandwidth cost \citep[\S 3.3.2]{deepseek-v3}.

The remaining DeepSeek-V3 choices follow from the same logic: master
weights and weight-gradient buffers stay in FP32; Adam first and
second moments are kept in BF16; the embedding matrix, output head,
MoE gating, normalization layers, and attention all stay in BF16 or
FP32 because their numeric properties (small dynamic range, sensitive
softmax, rare-token underflow risk) do not match what FP8's per-tile
scaling is designed to absorb \citep[\S 3.3.1]{deepseek-v3}. The
empirical envelope: relative loss error against a BF16 baseline stays
$\leq 0.25\%$ at scale \citep[\S 3.3]{deepseek-v3}, and the full
14.8T-token pre-training reported zero irrecoverable spikes
\citep[abstract]{deepseek-v3}. TorchTitan's per-tensor Float8 path is
considerably less aggressive but generic: the documented combined
speedup on Llama-3.1 with FP8 + FSDP2 + TP is reported in the
30--65\% range over BF16 \citep[\S 2.2.4]{torchtitan2024}, depending
on scale and which Linear layers are covered.

\begin{intuition}
Fine-grained scaling is what makes FP8 viable for full pre-training.
A single per-tensor scale must compromise between outliers (which need
a large $Z$) and bulk values (which need a small $Z$ for precision).
Per-tile scaling lets each $128$-element tile choose its own $Z$, so
outliers in one tile do not blow up precision in adjacent tiles.
Returning to math: \cref{eq:fp8-gemm} generalizes to per-tile
$Z_X^{(t)}$ and $Z_W^{(b)}$ along the inner dimension, and the GEMM
accumulates across tiles in CUDA-core FP32 every $N_C$ elements. The
effective dynamic range of an E4M3 tile with scale $Z$ is $Z \cdot
[-448, 448]$, so per-tile $Z$ recovers the dynamic range that NVIDIA's
hybrid was buying with E5M2's extra exponent bit, while E4M3's spare
mantissa bit becomes pure precision gain.
\end{intuition}

\subsection{Sub-FP8: NVFP4 and BitNet 1.58-bit}
\label{sec:mp-sub-fp8}

Below FP8 the field splits into two distinct programs.

\paragraph{NVFP4 (FP4 E2M1).} Blackwell-class GPUs (B100, B200) ship
native FP4 Tensor Cores with hardware MXScale: a per-block scaling
scheme operating on $32$-element blocks, where each block carries an
8-bit shared exponent. The format itself is E2M1 (2 exponent bits, 1
mantissa, dynamic range $\pm 6$ before the per-block scale)
\citep[\S 3.3, format zoo]{deepseek-v3}.\footnote{KB note:
\texttt{kb/notes/training/mixed-precision-and-stability.md\#1-3} and
\texttt{\#3-1}; format-zoo entry derived there from Blackwell white-
paper material.} The compute-throughput gain over FP8 is roughly
$2\times$ by spec, mirroring the FP16-to-FP8 step. The fine-grained-
plus-promotion recipe of \cref{sec:mp-fp8} extends naturally — replace
the $128$-element tile with the $32$-element MXScale block and lift
the same accumulator-promotion machinery — but whether it transfers
without quality loss is open.

\paragraph{BitNet 1.58.} \citet{ma2024-bitnet} push the precision
axis to its extreme: every weight is ternary in $\{-1, 0, +1\}$
(hence $\log_2 3 \approx 1.585$ bits per weight), and activations are
INT8. The construction replaces \texttt{nn.Linear} with
\texttt{BitLinear}, trained from scratch, using an \emph{absmean}
quantization function:
\begin{align}
\widetilde{W} &= \mathrm{RoundClip}\!\left(\tfrac{W}{\gamma + \epsilon}, -1, 1\right),
   \label{eq:bitnet-quant} \\
\mathrm{RoundClip}(x, a, b) &= \max(a, \min(b, \mathrm{round}(x))),
   \label{eq:roundclip} \\
\gamma &= \tfrac{1}{nm} \sum_{ij} |W_{ij}|,
   \label{eq:bitnet-gamma}
\end{align}
verbatim from \citet[\S 2, eqs.\ 1--3]{ma2024-bitnet}.\footnote{KB
excerpt: \texttt{kb/excerpts/ma2024-bitnet.md\#sec-2}.} Activations
are scaled per-token to $[-Q_b, Q_b]$ rather than $[0, Q_b]$ to
eliminate zero-point quantization \citep[\S 2]{ma2024-bitnet}. The
matrix multiply at inference reduces to integer addition (no floating-
point multiply) because the weight set is $\{-1, 0, +1\}$, saving
$\sim$71$\times$ arithmetic energy on 7nm chips by the
\citet{ma2024-bitnet} energy model \citep[\S 3]{ma2024-bitnet}.

The headline empirical claim is matched-perplexity at 3B
parameters: BitNet b1.58 at 3B reaches $9.91$ PPL versus LLaMA-LLM
3B at $10.04$, while running $2.71\times$ faster and using
$3.55\times$ less GPU memory \citep[\S 3, Table 1]{ma2024-bitnet}.
At $70$B the throughput claim is starker: $176$ vs.\ $16$ max batch
size, $2977$ vs.\ $333$ tokens/s \citep[\S 3, Table 3]{ma2024-bitnet}.
The architecture imports the LLaMA-alike stack — RMSNorm, SwiGLU,
RoPE, no biases — to enable drop-in inference under llama.cpp / vLLM
\citep[\S 2]{ma2024-bitnet}.

\subsection{Stability mitigations}
\label{sec:mp-stability}

Pre-training loss spikes — sudden, sometimes-irrecoverable jumps in
$\mathcal{L}$ — are observed across all frontier-scale runs (PaLM
\citep{chowdhery2022}, OPT, GLM-130B). The mitigations stack
across precision regimes:

\paragraph{Loss scaling.} Eq.~(\ref{eq:loss-scale}) for FP16; not
required for BF16; replaced by per-tile / per-block dynamic scales
for FP8 (Eq.~\ref{eq:fp8-gemm}).

\paragraph{Gradient clipping.} Clip $\|g\|_2 \leq c$ before the
optimizer step. The standard fix for gradient-norm explosion under LR
overshoot or batch-size transitions; orthogonal to precision.

\paragraph{Embedding LR scaling and rare-token handling.} Token
embeddings can grow unboundedly under weight decay if the optimizer's
preconditioning interacts poorly with embedding sparsity, and rare-
vocabulary entries' embedding columns receive small gradient
magnitudes that underflow in FP8. DeepSeek-V3's mitigation is direct:
keep embeddings in BF16/FP32 even when surrounding GEMMs are FP8
\citep[\S 3.3.1]{deepseek-v3}.

\paragraph{Attention-softmax kept higher-precision.} Softmax inputs
$QK^\top$ grow in magnitude as training proceeds, creating numeric
overflow when $|QK^\top|$ exceeds E4M3's $\pm 448$ bound. DeepSeek-V3
keeps the attention operator in BF16 for this reason \citep[\S
3.3.1]{deepseek-v3}; complementary fixes are \emph{QK-clip} (clamping
$QK^\top$ pre-softmax) and \emph{QK-LayerNorm} (normalizing $Q$ and
$K$ before the dot product, as adopted in LLaMA-2 and successors).

\paragraph{Optimizer-precision coupling.} Muon's update RMS,
$\sqrt{1/\max(A, B)}$ for an $A\times B$ matrix parameter, combined
with growing weight RMS, ``exceeds the high-precision range of
bf16'' without weight decay \citep[\S 2.2]{muon-moonlight2025}.\footnote{KB
excerpt: \texttt{kb/excerpts/muon-moonlight2025.md\#sec-2-2-rms} and
\texttt{\#sec-2-2-wd}.} The Muon-Moonlight fix is mandatory weight
decay plus per-shape $\sqrt{\max(A, B)} \cdot 0.2$ scaling that pulls
the update + weight back into BF16's representable range. Optimizer
choice and precision choice are not orthogonal; each new optimizer
needs a precision-pairing study.

\paragraph{MoE expert collapse.} The auxiliary load-balancing loss can
fail under FP8, especially for the expert-routing softmax. DeepSeek-
V3's auxiliary-loss-free load balancing addresses this in tandem with
the FP8 recipe \citep[abstract]{deepseek-v3}.

\paragraph{Communication precision.} The all-reduce of gradients in
distributed training is a separate precision choice.
\citet{shoeybi2019-megatron} showed that FP16 all-reduce with FP32
reduction on a master node maintains stability for tensor-parallel
runs. DeepSeek-V3's cross-node all-to-all kernels use FP8 dispatch and
BF16 combine \citep[\S 3.2.2]{deepseek-v3}, with the bandwidth budget
split between InfiniBand and NVLink.

\begin{contradiction}
NVFP4 / FP4 production status is not settled. Blackwell-class
hardware ships native FP4 Tensor Cores and several Blackwell-deploying
vendors are running FP4 pilots, but as of 2026-Q2 no major lab has
publicly confirmed FP4 \emph{pre-training} of a frontier model. The
\citet{scaling-laws-precision-2024} effective-parameter prediction
suggests substantial loss below $\sim$7 bits, but the prediction is
fitted on small models. BitNet 1.58-bit's matched-FP16 quality
\citep{ma2024-bitnet} is replicated independently only thinly past
$\sim$7B parameters, and the published BitNet b1.58 2B4T tech report
\citep{bitnet-2b4t-2025} does not yet have a frontier-scale companion.
DeepSeek-V3's FP8 result is the strongest production data point we
have at this precision level, and it is single-source: no other
publicly disclosed lab has matched the 14.8T-token zero-spike FP8
envelope \citep[abstract]{deepseek-v3}. Treat the entire sub-FP8
regime as open, and treat FP8 production stability as resolved
\emph{contingent on} the DeepSeek-V3 recipe carrying over to other
architectures.
\end{contradiction}

\begin{intuition}
Precision as compression. Each precision step (FP32 $\to$ BF16 $\to$
FP8 $\to$ FP4) halves storage and roughly doubles compute throughput
on supporting hardware. Training in lower precision is lossy
compression of the gradient computation, recovering quality through
(a) keeping a high-precision master copy that aggregates many small
updates and (b) per-block scaling that makes the loss tile-local
rather than tensor-wide. Returning to canonical form: the master-
weights pattern in Eq.~(\ref{eq:amp}) is exactly the lossy-
compression-with-error-correction pattern, and per-tile scaling in
Eq.~(\ref{eq:fp8-gemm}) is exactly tile-local quantization with
shared metadata. Each ingredient is independent; the combinations are
the recipes.
\end{intuition}

\subsection{Forward link}
\label{sec:mp-forward}

The mixed-precision and stability machinery of \cref{sec:mp-fp8} and
\cref{sec:mp-stability} concludes the pre-training scaffolding —
\cref{sec:mixed-precision} sits as the last of the upstream stages
that determine \emph{what base model} the post-training pipeline
inherits. From here we cross into post-training. The next section
(\cref{sec:sft}) walks supervised fine-tuning: the cross-entropy
objective with prompt-masking, the data-curation lineage from FLAN
through InstructGPT, Llama-3 rejection-sampling SFT, T\"ulu-3, and
the 2025 R1-Distill / s1 results that show reasoning capability
transfers through pure SFT. SFT inherits the mixed-precision recipe
of this section unchanged — typically BF16 working precision, FP32
master weights, and a learning rate roughly $10\times$ smaller than
pre-training peak — and the stability concerns shift from gradient-
underflow to data-mixture and prompt-mask convention.
